\documentclass[11pt]{article}
\usepackage[margin=1.1in]{geometry}
\usepackage[T1]{fontenc}
\usepackage{lmodern}
\usepackage{microtype}
\usepackage{natbib}
\usepackage{float}
\usepackage{amssymb}
\usepackage{../shared/smcbvs}
\bibliographystyle{abbrvnat}
\setcitestyle{authoryear,round}
\linespread{1.06}
\setlength{\emergencystretch}{2em}
\newcommand{\bridgepending}{The short bridging argument that puts the statement in this reader form awaits its final
independent check.}
% "Why it matters" line, and the list of assumptions with the place where each is judged.
\newcommand{\why}[1]{\par\noindent\textit{Why it matters.} #1\par\smallskip}
\newcommand{\asmlist}[1]{\par\smallskip\noindent{\small\textit{Assumptions, and where they are judged.} #1}\par\smallskip}
\newcommand{\SA}[1]{Assumptions~\S#1}

\title{What we prove: consistency of the SMC-BVS method}
\author{SMC-BVS project}
\date{Draft v2, \today}

\begin{document}
\maketitle

\begin{abstract}
\noindent
We prove five things about the SMC-BVS method, a sequential Monte Carlo method for Bayesian variable selection in the
linear model, with the prior exactly as the software sets it. The posterior finds the true model, and it finds the true
distribution of the response. The software, as coded, reports the true model. Its evidence estimate is unbiased on every
design, and, on designs without leakage and on the event that its calibration succeeds, its printed estimates are
consistent (no claim is made about how likely that event is). The rate of the first result cannot be
improved. The results are asymptotic in the sample size $n$ and allow many more predictors than observations.
The section ``Are the assumptions reasonable?'' (cited below as ``Assumptions~\S\,$x$'') judges every assumption used here.
\end{abstract}

% =====================================================================================================
\section{The results in one page}\label{sec:results}

The data are $n$ observations of a response and $p$ predictors, with $p$ possibly much larger than $n$; the response
follows a sparse Gaussian linear model. The method is a sequential Monte Carlo sampler for the posterior over models
under a fixed prior, followed by a local search; nothing in it is tuned by the user. We prove five results. Table~\ref{tab:results}
lists them; Sections~\ref{sec:model-result}--\ref{sec:estimates} state them, one section each, in the same format: what
the result says in words, why it matters, the statement, and its assumptions with the place in Section~A where each one is judged.

\begin{table}[ht]
\centering\footnotesize
\setlength{\tabcolsep}{3.5pt}
\renewcommand{\arraystretch}{1.12}
\caption{The five results. ``Checked'' means independently checked; statuses are those of the project's registry.}\label{tab:results}
\begin{tabular}{@{}>{\raggedright\arraybackslash}p{0.17\linewidth}>{\raggedright\arraybackslash}p{0.3\linewidth}>{\raggedright\arraybackslash}p{0.26\linewidth}>{\raggedright\arraybackslash}p{0.2\linewidth}@{}}
\toprule
Result & In words & Key assumptions & Status \\
\midrule
1.~The posterior finds the true model (Thms~\ref{thm:post-random}, \ref{thm:post-fixed}) &
The posterior mass outside the true model is of order $\rn$; the highest posterior and median probability models equal
the true model with probability tending to one. &
Gaussian linear model; $p\le n^A$; a fixed number of fixed nonzero coefficients; an eigenvalue floor on small sets of
predictors. &
Proved and checked; one bound in the fixed-design case rests on two supporting results not yet checked. \\
2.~It finds the true distribution (Thm~\ref{thm:dist}) &
The posterior contracts around the true conditional law of the response, at radius of order $\rn$ in the Hellinger distance. &
i.i.d.\ observations; sparse Gaussian linear model; light-tailed predictors; no near-exact collinearity; no lower bound on
the coefficients. &
Proved and checked. \\
3.~The method, as coded, reports the true model (Thm~\ref{thm:reported}) &
With probability tending to one, every run that produces output reports exactly the true model, whatever its seed. &
Gaussian rows; fixed signals; no true predictor masked by $15$ others; the predictors that can imitate the signal are few,
and local. &
Proved and checked; a revision of the statement is pending. \\
4.~The run's estimates (Prop.~\ref{prop:unbiased}, Thm~\ref{thm:outputs}) &
The evidence estimate is unbiased on every design. On designs without leakage the printed inclusion probabilities,
median model and evidence estimate are consistent. &
Unbiasedness: none on the design. Consistency: as Result~3, and no null predictor correlated with a true one. &
Per island, proved and checked; the pooled form is not yet checked. Consistency: proved and checked. \\
5.~The rate cannot be improved (Rem.~\ref{rem:sharp}) &
No design gives a rate faster than $n^{-1/2}$; on two families of designs the rate $\rn$ is attained. &
None for the first part; orthogonal or Gaussian null predictors for the second. &
Proved and checked. \\
\bottomrule
\end{tabular}
\end{table}

\paragraph{What ``asymptotic'' means.} The results are statements as $n\to\infty$, with $p$ allowed to grow polynomially
in $n$ and to exceed $n$. The constants in the proofs are explicit where they are displayed and conservative; they are
not part of any assumption we judge. One lower bound, Remark~\ref{rem:sharp}(a), holds at every $n\ge3$. At moderate $n$ and
large $p$ the posterior is not close to a point mass (Remark~\ref{rem:sharp}(c)); the theorems describe the limit.

\paragraph{What is not covered.} The parallel engine and runs with many islands are not covered by any theorem here; the
supplement (Section~\ref{sec:supp}) lists the further results we have, each with the reason it matters, and the open
questions.

% =====================================================================================================
\section{Setting, the method, and the shared assumptions}\label{sec:setting}

This section fixes the data, the model, the posterior and the few objects of the method that the results talk about.
The methodology section describes the method in full; Appendix~\ref{app:notation} lists the symbols.

\paragraph{Notation.} $\set{k}:=\{1,\dots,k\}$; $\norm{\cdot}$ is the Euclidean norm; $\lmin(\cdot)$ is the smallest
eigenvalue of a symmetric matrix; $\abs{S}$ is the size of a set; $\ind{E}$ is an indicator; $a\vee b:=\max(a,b)$. A star marks a true value,
and the superscript $\mathrm r$ marks raw data. $V_n=\OP(a_n)$ means that $\abs{V_n}/a_n$ is bounded in probability, and $V_n=\oP(a_n)$ that it tends to zero in
probability; \emph{with probability tending to one} means that the probability of the event tends to one. For positive
sequences, $a_n=O(b_n)$ (or $a_n\lesssim b_n$, $b_n\gtrsim a_n$) means $a_n\le Cb_n$ for a constant $C$, $a_n=o(b_n)$
means $a_n/b_n\to0$, and $a_n\gg b_n$ means $b_n=o(a_n)$. A quantity is \emph{fixed} if
it does not depend on $n$; all results are asymptotic in $n\to\infty$, along a sequence of problems in which $p$, the
design and the coefficients may change with $n$.

\subsection{Data, model and posterior}\label{sec:model}

We observe a response $\yr\in\R^n$ and $\praw$ predictors. The method removes every predictor that is constant in the
sample (its effect is absorbed by the intercept); $\Xr$ is the $n\times p$ matrix of the $p$ retained predictors. The data
follow a Gaussian linear model with an intercept,
\begin{equation}\label{eq:model}
  \yr=\alpha^*\one_n+\Xr\bstar+\varepsilon,\qquad \varepsilon\sim\Normal_n(0,\sigs^2I_n),
\end{equation}
with $\sigs^2>0$ and $\varepsilon$ independent of $\Xr$. A \emph{model} is a set $S\subseteq\set{p}$; the \emph{true
model} is $\Sz:=\{j:\beta^*_j\neq0\}$, of size $\sz$. Its members are the \emph{true predictors}; the others are
\emph{null predictors}. In a \emph{fixed design} $\Xr$ is deterministic; in a \emph{random Gaussian design} its rows are
i.i.d.\ $\Normal_p(m_n,\Sig)$, and $W=(W_1,\dots,W_p)\sim\Normal_p(0,\Sig)$ denotes a centred copy of one row, with
sub-vectors $W_S$ and sub-matrices $\Sigma_{n,S}$.

\paragraph{What the method sees.} The method centres and scales every predictor and the response: $X:=P^\perp_{\one}\Xr
D^{-1}$ and $y:=P^\perp_{\one}\yr/s_y$, with $P^\perp_{\one}$ the centring matrix, $D$ the diagonal matrix of column sample
standard deviations and $s_y$ that of the response. Then $P^\perp_{\one}\yr=X\tstar+P^\perp_{\one}\varepsilon$ with the
\emph{standardised coefficients} $\tstar:=D\bstar$ (same support $\Sz$), the \emph{mean signal} $\mu:=X\tstar$, and
$\tmin:=\min_{j\in\Sz}\abs{\theta^*_j}$. For a model $S$, $X_S$ holds its columns, $G_S:=X_S^\top X_S$, $\Rhat_S:=G_S/(n-1)$
is its sample correlation matrix and $P_S$ the projection onto its column span.

\paragraph{The prior and the posterior, as coded.} The method fits a working Gaussian linear model,
$y\mid S,\theta_S,\sigma^2\sim\Normal_n(X_S\theta_S,\sigma^2I_n)$, to the standardised data; the theorems are proved
under the true model~\eqref{eq:model}, not under this working model. It considers the models of at most $K:=\min(p,\lfloor
n/2\rfloor)$ predictors, $\Om:=\{S:\abs{S}\le K\}$, and puts on them the beta-binomial prior $\BB(1,p)$ \citep{ScottBerger2010},
restricted to $\Om$; given $S$, a $g$-prior with $g=n$ and a unit ridge on the coefficients \citep{Zellner1986}; and an
inverse-gamma prior $\IG(0.01,0.01)$ on the noise variance. Integrating out the coefficients and the variance gives each
model a marginal likelihood $m(S)$ in closed form (methodology section, eq.~(6)), and the \emph{posterior over models}
\begin{equation}\label{eq:posterior}
  \pis{S}=\frac{e^{F(S)}}{Z},\qquad F(S):=\log\{\text{prior of }S\}+\log m(S),\qquad Z:=\sum_{S'\in\Om}e^{F(S')} .
\end{equation}
$F(S)$ is the \emph{log score} of $S$, and the normalising constant $Z$ is the \emph{evidence}. Adding one predictor to a model of size $k$ multiplies the prior by
$(k+1)/(2p-1-k)$, and, in a well-conditioned model, multiplies $m(S)$ by about $(n+1)^{-1/2}$ (an Occam factor) times a
gain of fit.

\paragraph{Posterior summaries.} The \emph{inclusion probability} of predictor $j$ is $\PIP_j:=\sum_{S\ni j}\pis{S}$;
the \emph{highest posterior model} is $\HPM:=\argmax_S\pis{S}$; the \emph{median probability model} is
$\MPM:=\{j:\PIP_j>1/2\}$ \citep{BarbieriBerger2004}; and the \emph{posterior miss} is $1-\pis{\Sz}$.

\begin{lemma}[The posterior miss controls the summaries]\label{lem:summaries}
For every data set, $\max_{j\in\set{p}}\abs{\PIP_j-\ind{j\in\Sz}}\le1-\pis{\Sz}$. If $\pis{\Sz}>1/2$, then
$\HPM=\MPM=\Sz$.
\end{lemma}

The proof is two lines: every model that misses a true predictor, or contains a null one, differs from $\Sz$.

\subsection{A run of the method, and what it returns}\label{sec:method}

A \emph{run} is one call of the software on a data set, with a seed and a computing allocation; a data set is
\emph{valid} when $n\ge3$, at least one predictor is given and every entry is finite. When $\Om$ is small enough the
method enumerates it; otherwise it runs the sampler. A run either \emph{produces output}, with status ``exact''
(enumeration) or ``completed'' (sampler; a \emph{completed run}), or stops without estimates. Four objects of a run appear
in the results.
\begin{itemize}[leftmargin=1.5em,itemsep=1pt]
  \item \emph{The start law.} The method partitions the predictors into blocks of at most $16$, by clustering on sample
    correlations, and scores every subset of each block exactly. The \emph{start law} is the posterior with all
    interactions between blocks $C_1,\dots,C_J$ ignored,
    $\pi_0(S)\propto w(\abs{S})\prod_{b=1}^{J}m(S\cap C_b)/m(\emptyset)$, where $w(\abs{S})$ is the prior weight of a
    model of that size; it can be sampled exactly. A maximiser of $\pi_0$ is a \emph{start mode} $\tSz$.
  \item \emph{The sampler.} A population of models moves from $\pi_0$ to the posterior through intermediate
    distributions, with temperatures and effort chosen by an internal \emph{calibration} that is frozen before the final
    run. The calibration \emph{succeeds} when it produces a sequence of temperatures (by its main construction or a
    fallback); otherwise the run stops. We write $\Acal$ for the event that the method uses the sampler and its
    calibration succeeds. The final run consists of four independent copies, the \emph{islands}.
  \item \emph{The reported model.} After the sampler, two local searches run (from the best model found, and from the
    start mode); each moves to the best neighbouring model (one predictor added, deleted or swapped) while that strictly
    improves the log score, for at most a capped number of moves ($\ell_n:=\lceil\log_2(n+p+1)\rceil$ for the first search,
    $\ell_n+\abs{\tSz}$ for the search from the start mode), and then scans the neighbours once more.
    The \emph{reported model} $\Shat$ is the best-scoring model the run has recorded. The method \emph{certifies}
    $\Shat$ (sets its local-optimality flag) when it enumerated $\Om$, or when $\Shat$ ends a local search whose last scan
    found no strict improvement; a search that reaches its cap and still finds an improving neighbour stops uncertified.
  \item \emph{The estimates.} Island $r$ estimates the inclusion probabilities and the evidence; the method pools the
    four islands with equal weights into $\PIPhat_j$ and $\Zhat$, and reports $\MPMhat:=\{j:\PIPhat_j>1/2\}$.
\end{itemize}

\subsection{Shared assumptions and conventions}\label{sec:standing}

The assumptions below are used across Sections~\ref{sec:model-result}--\ref{sec:estimates}; the others are stated
where they are first used. Our simulation designs (S1--S6, the strong-factor study, the designs IND, AR and BLK of the
planned runs at $p=10^6$, and FAC, a contingent design planned at $p=10^4$ only) are described in Assumptions~\S2: S1 and IND have independent predictors; S2--S4,
S6, AR (autoregressive) and BLK (blocks) are locally correlated; S5, the strong-factor study and FAC have pervasive common
factors.

Two words recur in Results~3 and~4. A true predictor is \emph{visible} if its signal does not cancel against any $15$ other
predictors (the method's blocks hold at most $16$); a null predictor \emph{leaks} if, given up to $15$ others, it can imitate
part of a true predictor's signal. Sections~\ref{sec:reported} and~\ref{sec:estimates} make both precise.
The assumptions section judges each of them and compares it with \citet{SongLiang2015}.

\begin{assumption}{GL}{Gaussian linear model}
The data follow~\eqref{eq:model}, with a fixed noise variance $\sigs^2>0$ and noise independent of the design. \hfill [\SA{3.1}]
\end{assumption}

\begin{assumption}{DIM}{Dimension}
$2\le p\le n^A$ for a fixed $A\ge1$. \hfill [\SA{3.1}]
\end{assumption}

\begin{assumption}{SP}{Sparsity}
The true model size $\sz\ge1$ is fixed. \hfill [\SA{3.1}]
\end{assumption}

\begin{assumption}{GR}{Gaussian rows}
The rows of $\Xr$ are i.i.d.\ $\Normal_p(m_n,\Sig)$, where $\Sig$ is a correlation matrix. \hfill [\SA{3.2}]
\end{assumption}

\begin{assumption}{EFS}{Population eigenvalue floor}
$\lmin(\Sigma_{n,S})\ge\varphi^2$ for every $S$ with $\abs{S}\le5(A+1)(\sz+1)\log n$, for a fixed $\varphi\in(0,1]$ (for
instance, $\lmin(\Sig)\ge\varphi^2$). \hfill [\SA{3.2}]
\end{assumption}

The conventions fix what ``the method'' means and idealise the computer; they are not statistical assumptions.
\emph{Convention \as{C1}} (the method): exact arithmetic, the default configuration, valid inputs and seed; a run that
meets it is \emph{admissible}, a sample size $n$ is admissible if some data set of that size is valid, and statements
hold for every valid seed and every outcome of the calibration.
\phantomsection\label{asm:C1}%
\emph{Convention \as{C2}} (delivery): an idealised machine that never overflows or fails, ideal random numbers, and an
allocation of memory and time no smaller than the method's own polynomial requirement formulas ask for; \emph{\as{C2p}} also lifts two fixed integer limits of the code (Appendix~\ref{app:conventions}).\phantomsection\label{asm:C2}\phantomsection\label{asm:C2p}
\emph{Convention \as{C3}} (sampler): exact arithmetic, ideal random numbers independent of the data and the calibration,
an exact start sampler, and Markov moves that do what their definitions say.\phantomsection\label{asm:C3}
Appendix~\ref{app:conventions} gives the full lists.

% =====================================================================================================
\section{Result 1: the posterior finds the true model}\label{sec:model-result}

\inwords{If the true model is small and its coefficients are fixed and nonzero, and no small group of predictors is
collinear, then the posterior puts all but a fraction of order $\rn$ of its mass on the true model. The highest posterior
model and the median probability model then equal the true model with probability tending to one, and every inclusion
probability is within order $\rn$ of $1$ for a true predictor and of $0$ for a null one. We state this first for Gaussian
rows, the case of our simulation designs, and then for a fixed design, with conditions on the sample.}

\why{It is the basic guarantee of a Bayesian selection method, here for the prior exactly as the software sets it: with
enough data the posterior points to the right predictors, and its inclusion probabilities can be read as near-certain
inclusion or exclusion.}

\begin{assumption}{FSb}{Fixed signals, raw scale}
$b_{\mathrm{lo}}\le\abs{\beta^*_j}\le b_{\mathrm{hi}}$ for every $j\in\Sz$, with fixed $0<b_{\mathrm{lo}}\le b_{\mathrm{hi}}<\infty$.
\end{assumption}

\begin{theorem}[Posterior concentration, random Gaussian design]\label{thm:post-random}
Assume \as{GL}, \as{DIM}, \as{SP}, \as{GR}, \as{FSb} and \as{EFS}. Then, under the joint law of the design and the noise,
\begin{gather}
  1-\pis{\Sz}=\OP\Big(\sqrt{\frac{\log n}{n}}\Big),\qquad \P\big(\HPM=\MPM=\Sz\big)\to1,\label{eq:post-random}\\
  \max_{j\in\set{p}}\big|\PIP_j-\ind{j\in\Sz}\big|=\OP\Big(\sqrt{\frac{\log n}{n}}\Big).\nonumber
\end{gather}
\end{theorem}

\asmlist{(1)~\as{GL} Gaussian linear model; (2)~\as{DIM} $p\le n^A$; (3)~\as{SP} fixed $\sz$ [\SA{3.1}]; (4)~\as{GR}
Gaussian rows; (5)~\as{FSb} fixed nonzero coefficients; (6)~\as{EFS} population eigenvalue floor on sets of
$O(\log n)$ predictors [\SA{3.2}].}

\status{\textbf{Proved and independently checked}: the fixed-design result it rests on (Theorem~\ref{thm:post-fixed},
in-probability form) and the transfer from population to sample conditions are each proved and independently checked.
Registry: Appendix~\ref{app:status}, row 1.2.}

No upper bound on the eigenvalues of $\Sig$ is needed, so strong common factors are allowed, and $p\gg n$ is allowed. In
primitive terms: a Gaussian linear model with i.i.d.\ Gaussian rows whose correlation matrix has smallest eigenvalue bounded below, a fixed number of
fixed nonzero coefficients and $p\le n^A$ suffice.

\paragraph{Fixed design.} The same conclusions hold for a deterministic design under conditions on the sample. They
use four more assumptions.

\begin{assumption}{BC}{Bounded coefficients}
$\max_{j}\abs{\theta^*_j}\le\bar\theta$ for a fixed $\bar\theta<\infty$. \hfill [\SA{3.2}]
\end{assumption}

\begin{assumption}{EF}{Eigenvalue floor on sets of bounded size}
$\lmin(\Rhat_S)\ge\varphi_0^2$ for every model $S$ with $\abs{S}\le2(A+1)\sz+2$, for a fixed $\varphi_0\in(0,1]$. \hfill [\SA{3.2}]
\end{assumption}

\begin{assumption}{NC}{No near-exact collinearity among sets of logarithmic size}
$\lmin(G_S)\ge g_0$ for every model $S$ with $\abs{S}\le3(A+1)(\sz+1)\log n$, for a fixed $g_0>0$; on the correlation
scale, a floor $g_0/(n-1)$ that vanishes as $n$ grows. \hfill [\SA{3.2}]
\end{assumption}

\begin{assumption}{BM}{Beta-min}
Either (a) $\tmin\ge C_\beta\,\sigs\rn$ with a fixed $C_\beta>0$ such that $C_\beta^2\varphi_0^2>c_0(\sz,A)$, an explicit
constant (Appendix~\ref{app:constants}); or (b) $\tmin\sqrt{n/\log n}\to\infty$. \hfill [\SA{3.2}]
\end{assumption}

\begin{theorem}[Posterior concentration, fixed design]\label{thm:post-fixed}
Let the design be fixed and assume \as{GL}, \as{DIM}, \as{SP}, \as{BC}, \as{EF}, \as{NC} and \as{BM}. Then, along every
sequence that satisfies these assumptions,
\begin{equation}\label{eq:post-op}
  \max_{j\in\set{p}}\big|\PIP_j-\ind{j\in\Sz}\big|\ \le\ 1-\pis{\Sz}\ =\ \OP\Big(\sqrt{\frac{\log n}{n}}\Big),\qquad
  \P\big(\HPM=\MPM=\Sz\big)\to1 .
\end{equation}
Moreover, there are $n_0$ and $C_{\mathrm E}$, depending only on $(\sz,A,\sigs^2,\bar\theta,\varphi_0,g_0,C_\beta)$, such that for
every $n\ge n_0$ and every design and $\tstar$ that satisfy the assumptions at that $n$,
\begin{equation}\label{eq:post-exp}
  \E\big[1-\pis{\Sz}\big]\ \le\ C_{\mathrm E}\sqrt{\frac{\log n}{n}},\qquad \P\big(\HPM=\MPM=\Sz\big)\ \ge\ 1-4C_{\mathrm E}\sqrt{\frac{\log n}{n}} .
\end{equation}
The expectation is over the noise, with the design fixed. Under form (b) of \as{BM}, $n_0$ also depends on how fast
$\tmin\sqrt{n/\log n}$ diverges.
\end{theorem}

\asmlist{(1)--(3)~\as{GL}, \as{DIM}, \as{SP} [\SA{3.1}]; (4)~\as{BC} bounded coefficients; (5)~\as{EF} eigenvalue floor
on sets of bounded size; (6)~\as{NC} no near-exact collinearity; (7)~\as{BM} beta-min [\SA{3.2}].}

\status{The in-probability conclusions~\eqref{eq:post-op}: \textbf{proved and independently checked}. The expectation
bound~\eqref{eq:post-exp}: \textbf{proved and checked, conditional on one supporting result} (a proposition giving a lower
bound on the posterior miss) \textbf{that is proved but not yet independently checked}; its other input, a lemma bounding
the tail of the posterior model size, is now independently checked. Registry: Appendix~\ref{app:status}, row 1.1.}

\subsection{Result 5: the rate cannot be improved}\label{sec:sharp}

\inwords{On every design, no model of at most $\min(K-1,p/2)$ predictors has posterior mass closer to one than about
$(\abs{S}+1)/(3\sqrt n)$; on two families of designs the expected posterior miss is at least of order $\rn$. At simulated sizes the posterior spreads over the true model and
its one-predictor extensions.}

\why{It shows that the rate of Result~1 is the right one, and why the posterior is not a point mass at simulated sizes
such as $n=200$.}

\begin{remark}[The rate cannot be improved, and what it means at moderate $n$]\label{rem:sharp}
\textup{(a)} \emph{Every design.} For every $n\ge3$, every standardised design with $p\ge2$, every response and every
model $S$ with $\abs{S}\le\min(K-1,p/2)$,
\begin{equation}\label{eq:sharp-a}
  \pis{S}\ \le\ \Big(1+\frac{\abs{S}+1}{3\sqrt n}\Big)^{-1}.
\end{equation}
So no model of that size, the true one included, can carry posterior mass closer to one than about
$(\abs{S}+1)/(3\sqrt n)$. For true models with $\sz\le\min(K-1,p/2)$, no design gives a rate faster than
$n^{-1/2}$, and a posterior miss of order $\rn$ is possible only if $\sz=O(\sqrt{\log n})$. The reason is that the prior and the Occam factor keep every one-predictor extension of a model
alive.

\textup{(b)} \emph{Witness designs.} Assume \as{GL}, that $\sz$ is fixed, the energy bounds $\norm{\mu}^2\le nM_\mu$ and
$\norm{\tstar}^2\le nM_\theta$ for fixed $M_\mu$, $M_\theta$, and that the null predictors are either orthogonal (to each
other and to the true ones) or have i.i.d.\ $\Normal(0,1)$ raw entries independent of the true predictors and of the noise, with $p\le n^A$.
Then, for every $p>\sz$, bounded or growing,
where $\pi=3.14\ldots$ (the number, not the posterior) in the next display,
\begin{equation}\label{eq:sharp-b}
  \E\big[1-\pis{\Sz}\big]\ \ge\ (1-o(1))\,(\sz+1)\,\frac{2(p-\sz)}{2p-1-\sz}\,
  \sqrt{\frac{\log\!\big(p\sqrt n/(\sz+1)\big)}{\pi n}},
\end{equation}
and, if moreover $p\to\infty$, $1-\pis{\Sz}\ge(1-\oP(1))\,(\sz+1)\sqrt{\log p/(\pi n)}$. On fully
orthogonal designs ($\sz<p\le n-1$) with fixed nonzero coefficients, the expectation bound~\eqref{eq:post-exp} applies
too, so there $\E[1-\pis{\Sz}]$ is of exact order $\rn$.

\textup{(c)} \emph{At moderate $n$} (a heuristic, not a theorem). The posterior miss is carried mainly by the models
$\Sz\cup\{j\}$ that add one null predictor. Each has posterior odds against $\Sz$ of about
$\frac{\sz+1}{2p}\,n^{-1/2}e^{T_j^2/2}$ (prior, Occam factor and fit), with $T_j^2$ approximately $\chi^2_1$. The factor
$e^{T_j^2/2}$ has no finite mean; over $p$ null predictors the largest $T_j^2$ is about $2\log p$, and the $p$ factors
average about $2\sqrt{\log p/\pi}$. Together these models therefore carry about $(\sz+1)\sqrt{\log p/(\pi n)}$ times the
weight of $\Sz$: the expression of~(b), here as a ratio to the weight of $\Sz$ rather than a bound on the posterior miss. With $\sz=5$, $n=200$ and $p=10^4$ this is about $0.7$. In single
synthetic draws (AR(0.6), $\sz=5$) the measured total weight of these models, relative to that of $\Sz$, was $1.0$--$1.2$
at $n=200$ and $p=10^4$--$10^6$, and $0.52$ at $n=400$, $p=10^4$. So at simulated sizes the posterior is spread over $\Sz$
and many one-predictor extensions of it.
\end{remark}

\asmlist{(a): none. (b): \as{GL}, a fixed $\sz$ [\SA{3.1}], and the witness designs, which describe where the rate is
attained and are not conditions of consistency.}

\status{(a) and (b): \textbf{proved and independently checked}; the exact-order statement after (b) uses the
expectation bound~\eqref{eq:post-exp} and has its status. Registry: Appendix~\ref{app:status},
row 1.3. (c): a heuristic computed by the Theory department, with one-draw measurements from the engineering
experiments; no registry row.}

% =====================================================================================================
\section{Result 2: the posterior finds the true distribution}\label{sec:distribution}

\inwords{Suppose the observations are i.i.d., the response follows a sparse Gaussian linear model, the predictors have light
tails, no small group of predictors is almost exactly collinear, and no single predictor carries an unbounded signal.
Then, for every large enough level $M$, with probability at least $1-n^{-c_{\mathrm H}M}$ the posterior puts mass at
most $n^{-c_{\mathrm H}M}$ on conditional laws of the response farther than $\sqrt{M\log n/n}$ from the true one, in the
Hellinger distance. No
lower bound on the coefficients is needed.}

\why{The posterior describes the response correctly even when some true predictors are too weak to be selected, so
prediction does not depend on perfect selection.}

\paragraph{Setting.} In this section only, the observations are $n$ i.i.d.\ copies $(x^{(i)},y^{(i)})$ of a pair $(x,y)$,
with $x\in\R^{\praw}$ the vector of all observed predictors; $\Dn$ is the data set, $\P^*$ its law, and the method runs on
the predictors that are not constant in the sample. The posterior is not defined if fewer than two predictors remain or
the response is constant in the sample; such data sets are counted in the exceptional event of the theorem. A posterior draw $(S,\beta,\sigma^2)$, mapped back to the original
units, defines the conditional law $\Normal(\bar y+(x-\bar x)^\top\beta,\sigma^2)$ of $y$ given $x$; the law of $x$ is not
estimated. Let $f$ be the joint density of $(x,y)$ with $x$ from its true law and $y\mid x$ from this fitted law, and $f^*$
the true joint density; then $d_H(f,f^*):=(1-\int\sqrt{ff^*})^{1/2}$ is the Hellinger distance between the fitted and the true conditional
laws, averaged over a new $x$. $\Pi(\cdot\mid\Dn)$ is the posterior of $(S,\beta,\sigma^2)$, and
$\zeta_j:=(x_j-\E x_j)/\sd(x_j)$. Fix a \emph{level} $M\ge1$; the \emph{window} $q_n:=\min\{\praw,\lceil C_qM\log
n\rceil\}$, $C_q:=2(s(A+2)+4)/\log2$, is the largest number of predictors the conditions constrain jointly.

\begin{localassumption}[Distribution-level setting]\label{asm:dist}
The integer $s\ge0$ and the constants $A\ge1$, $L\ge1$, $c_\psi>0$ and $\bar b>0$ are fixed.
\begin{asmitems}
  \item\label{asm:dist-model} \emph{Gaussian linear model for i.i.d.\ pairs.} $y=\alpha^*+x^\top\bstar+\varepsilon$ with
    $\varepsilon\sim\Normal(0,\sigs^2)$, $\sigs^2>0$, independent of $x$; $\alpha^*$, $\bstar$, $\sigs$ and the law of $x$ may
    depend on $n$.
  \item\label{asm:dist-sparse} \emph{Sparsity and dimension.} At most $s$ coordinates of $\bstar$ are nonzero, and
    $2\le\praw\le n^A$.
  \item\label{asm:dist-subg} \emph{Sub-Gaussian predictors on sparse directions.} Every predictor has positive variance,
    and $\E\exp(\lambda u^\top\zeta)\le\exp\big(L^2\lambda^2\Var(u^\top\zeta)/2\big)$ for every $\lambda\in\R$ and every
    $u\in\R^{\praw}$ with at most $q_n$ nonzero entries.
  \item\label{asm:dist-nec} \emph{No near-exact collinearity.} $\lmin(\Corr(x_T))\ge c_\psi/n$ for every set $T$ of at
    most $q_n$ predictors.
  \item\label{asm:dist-bs} \emph{Bounded signal.} $\abs{\beta^*_j}\sd(x_j)\le\bar b\,\sigs$ for every $j$.
\end{asmitems}
\end{localassumption}

\begin{theorem}[Distribution-level consistency]\label{thm:dist}
There are constants $c_{\mathrm H}>0$, $M_0$ and $n_0$, depending only on $(s,A,L,c_\psi,\bar b)$, such that the following
holds. Fix a level $M\ge M_0$ and assume Assumption~\ref{asm:dist}, with the window $q_n$ of this $M$. Then for every
$n\ge n_0$ with $\sqrt n/\log n\ge M$,
\begin{equation}\label{eq:dist-i}
  \P^*\Big[\,\Pi\Big(d_H(f,f^*)>\sqrt{M\log n/n}\ \Big|\ \Dn\Big)>n^{-c_{\mathrm H} M}\Big]\ \le\ n^{-c_{\mathrm H} M}.
\end{equation}
\end{theorem}

\asmlist{(1)~Assumption~\ref{asm:dist}\ref{asm:dist-model}: Gaussian linear model [\SA{3.1}]; (2)~\ref{asm:dist-sparse}:
sparsity and $\praw\le n^A$ [\SA{3.1}]; (3)~\ref{asm:dist-subg}: light-tailed predictors [\SA{3.2}, law of the rows];
(4)~\ref{asm:dist-nec}: no near-exact collinearity [\SA{3.2}]; (5)~\ref{asm:dist-bs}: bounded signal [\SA{3.2}, signal
size].}

\status{\textbf{Proved and independently checked}. Registry: Appendix~\ref{app:status}, row 2.1.}

\begin{remark}[On the assumptions]\label{rem:dist-asm}
No beta-min condition and no fixed lower bound on eigenvalues are assumed. Jointly Gaussian predictors satisfy Assumption~\ref{asm:dist}\ref{asm:dist-subg} with $L=1$, and every correlation matrix whose $q_n\times q_n$
principal blocks have smallest eigenvalue at least a fixed $\psi>0$ satisfies Assumption~\ref{asm:dist}\ref{asm:dist-nec}
for all large $n$. Condition~\ref{asm:dist-nec} keeps the Occam factor of each extra predictor in force; without it the
radius grows by one factor $\sqrt{\log(\praw\vee n)}$ (Proposition~\ref{prop:dist}). The same statement holds for every
level $M$ up to $\sqrt n/\log n$, with the same constants, It also carries over to the list of models that a completed run returns,
when that list holds posterior mass at least $n^{-c_{\mathrm H}M/2}$, a condition the run does not check
(Corollary~\ref{cor:dist-output}).
\end{remark}

% =====================================================================================================
\section{Result 3: the method, as coded, reports the true model}\label{sec:reported}

\inwords{If every true predictor is visible to a block-by-block search, the predictors that can imitate the signal form a
set $Q$ of at most $n^b$ predictors with $b<\tfrac12$, and regressions within $Q$ are local, then there is an event,
determined by the data alone and of probability tending to one, on which every run that produces output reports exactly
the true model, whatever its seed and the outcome of its calibration. The start mode contains every true predictor and
possibly some members of $Q$; the local search from the start mode removes the spurious ones, one per move.}

\why{It is a guarantee about what the software prints, at the effort its own calibration chooses and for every seed, not only
about the posterior it targets.}

\paragraph{How the signal looks to a block-by-block method.} The start law sees interactions within blocks of at most
$16$ predictors only, so the conditions of this section are about partial covariances given at most $15$ other
predictors. With $W\sim\Normal_p(0,\Sig)$, $\Cov(W_j,W_l\mid W_U)$ is the partial covariance given $(W_k)_{k\in U}$, and
$\bar\beta_{U,l}:=\Sigma_{n,UU}^{-1}\Sigma_{n,Ul}$ is the population regression coefficient of $W_l$ on $W_U$, padded
with zeros to a vector indexed by a larger set when needed. A
\emph{leakage set} is a deterministic $Q$ with $\Sz\subseteq Q\subseteq\set{p}$ that contains every predictor able to
imitate part of the signal; its \emph{lattice} $\Lat(Q):=\{T:\Sz\subseteq T\subseteq Q\}$ is the set of models between the
true model and $Q$. The signal is fixed on the scale the method sees:

\begin{assumption}{FSt}{Fixed signals, standardised scale}
$\yr=\alpha\one+X\tstar+\varepsilon$ with $\varepsilon\sim\Normal_n(0,\sigs^2I_n)$, a fixed $\sigs^2>0$, noise independent
of $\Xr$, a deterministic $\tstar$, $\sz\ge1$ fixed, and $\theta_{\mathrm{lo}}\le\abs{\theta^*_j}\le\theta_{\mathrm{hi}}$ for
$j\in\Sz$, with fixed $0<\theta_{\mathrm{lo}}\le\theta_{\mathrm{hi}}$. \hfill [\SA{3.2}]
\end{assumption}

\begin{assumption}{V}{Visibility}
For every $j\in\Sz$ and every $U\not\ni j$ with $\abs{U}\le15$, $\abs{\Cov(W_j,W^\top\tstar\mid W_U)}\ge a$, for a fixed
$a>0$. \hfill [\SA{3.3}]
\end{assumption}

In \as{FSt} the raw coefficients $\bstar=D^{-1}\tstar$ depend on the design; \as{FSt} takes the place of \as{GL} and \as{SP}
in Results~3 and~4.

\begin{example}[Visibility and leakage with three predictors]\label{ex:visibility}
Let $\Sz=\{1,2\}$ with $\tstar_{\Sz}=(1,-1)$, $\Corr(W_1,W_2)=\varrho$, and let $W_3$ be a null predictor. With $U=\emptyset$,
$\Cov(W_1,W^\top\tstar)=1-\varrho$, so \as{V} requires $1-\varrho\ge a$: as $\varrho\to1$ the two true predictors cancel and
each becomes invisible. If $\Corr(W_3,W_1)$ is fixed and nonzero, $W_3$ can imitate part of the signal of $W_1$, so it
must belong to the leakage set $Q$.
\end{example}

Fix $b\in(0,\tfrac12)$ and $b_H>0$ with $15b_H+b<\tfrac12$.

\begin{localassumption}[Polynomial leakage set with local regressions]\label{asm:reported}
\begin{asmitems}
  \item\label{asm:rr-leak} \emph{Leakage set.} There is a deterministic $Q\supseteq\Sz$ with $\abs{Q}\le n^b$ and, for
    large $n$, $\abs{Q}\le K$ and ($Q=\set{p}$ or $K\ge2\abs{Q}+32\sz+1$), such that: (i) \emph{no leak outside $Q$:}
    there is a deterministic $\varpi_n$ with $\varpi_n\sqrt{n/\log n}\to0$ such that
    $\abs{\Cov(W_l,W_j\mid W_U)}\le\varpi_n$ for all $l\notin Q$, $j\in\Sz\setminus U$ and $\abs{U}\le15$; and (ii)
    \emph{locality:} there is a deterministic $r_n\to0$, and every $l\notin Q$ has a deterministic
    $\Hl_l\supseteq\{l\}\cup\Sz$ with $\abs{\Hl_l}\le n^{b_H}$, such that for all $l\notin Q$, all $U\not\ni l$ with
    $\abs{U}\le15$ and $U_0:=U\cap\Hl_l$: $\abs{\Cov(W_l,W_v\mid W_{U_0})}\le r_n$ and
    $\abs{\Cov(W_j,W_v\mid W_{U_0})}\le r_n$ for all $v\in U\setminus\Hl_l$ and $j\in\Sz\setminus U_0$.
  \item\label{asm:rr-floor} \emph{Floor on the leakage set.} $\lmin(\Sigma_{n,Q\cup\{l\}})\ge\varphi^2$ for every
    $l\in\set{p}$, with $\varphi$ as in \as{EFS}.
  \item\label{asm:rr-local} \emph{Local regressions within the leakage set.} There are deterministic sequences
    $d_n=o(\log n)$ and $L_n$, and for every $l\in\set{p}$ a deterministic $N_l\subseteq Q\setminus\{l\}$ with
    $\abs{N_l}\le d_n$, such that $\norm{\bar\beta_{T,l}-\bar\beta_{T\cap N_l,l}}_1\le L_n$ for every $T\in\Lat(Q)$ with
    $l\notin T$, where $L_n\to0$ and $L_n^2\log(\abs{Q}+1)=o(\log n)$.
\end{asmitems}
\end{localassumption}

In~(a), $32\sz=2\cdot16\sz$, with $16$ the method's block-size cap: with probability tending to one the start mode has at
most $2(\abs{Q}+16\sz)$ predictors, so this condition leaves room to add a missing true predictor to it.

\begin{theorem}[The reported model equals the true model]\label{thm:reported}
Let $\Sig$ be positive definite and assume \as{GR}, \as{DIM}, \as{FSt}, \as{V}, \as{EFS}, Assumption~\ref{asm:reported}
and Convention~\as{C1}. Then, under the joint law of the design and the noise, there are data events $\mathcal G_n$ with
$\P(\mathcal G_n)\to1$ on which \emph{every} admissible run that produces output reports $\Shat=\Sz$, whatever its seed and the
outcome of its internal calibration. In particular,
\begin{equation}\label{eq:reported}
  \P(\text{the run produces output and }\Shat\neq\Sz)\to0 .
\end{equation}
If moreover Convention~\as{C2p} holds, then
\begin{equation}\label{eq:delivery}
  \P(\text{the run produces output})=1\ \ \text{for every admissible }n\ge3,\qquad \P(\Shat=\Sz)\to1 .
\end{equation}
\end{theorem}

Convention~\as{C2p} lifts two integer limits of the code; in the default configuration neither can be reached on any input
R can represent ($n,\praw\le2^{31}-1$), so this matters only for the form as $n\to\infty$ (Appendix~\ref{app:conventions}).

\asmlist{(1)~\as{GR} Gaussian rows; (2)~\as{DIM} [\SA{3.1}]; (3)~\as{FSt} fixed signals [\SA{3.2}]; (4)~\as{V}
visibility [\SA{3.3}]; (5)~\as{EFS} population floor [\SA{3.2}]; (6)~Assumption~\ref{asm:reported}(a) polynomial
leakage set, and (b)--(c) floor and local regressions in it [\SA{3.3}]; (7)~Conventions \as{C1}, and \as{C2p} for delivery
(Appendix~\ref{app:conventions}).}

\status{\textbf{Proved and independently checked}, with no open condition: the recorded master is a strengthening of
this statement, and the two code-reading audits that its delivery route takes as hypotheses are independently checked.
The statement that the integer limits of \as{C2p} cannot be reached on any representable input: \textbf{proved and
independently checked}. Registry: Appendix~\ref{app:status}, row 3.1.}

\begin{remark}[On the assumptions]\label{rem:reported}
(i) \emph{Why $b<\tfrac12$.} At model size about $\sqrt n$ the prior-plus-Occam penalty of one more predictor
($\approx\log p+\tfrac12\log n$) stops dominating the best of $p$ noise gains ($\approx\log p$); for larger $b$ the key events of the argument are
expected to fail (a heuristic; whether the conclusion fails is open).
(ii) Designs that satisfy the assumptions include correlation that is local in a metric and bounded blocks; of our
designs, the independent and locally correlated ones are inside (\SA{3.3}). The assumptions fail for pervasive common factors when
$p\gg\sqrt n$, because every null predictor then imitates part of the signal. There Results~1 and~2 and the
unbiasedness of Result~4 still hold, and certified screening (Theorem~\ref{thm:screen}), which has no leakage
condition, covers these designs when its own conditions hold.
\end{remark}

% =====================================================================================================
\section{Result 4: the run's estimates}\label{sec:estimates}

\inwords{The method prints an evidence estimate, pooled inclusion probabilities and the median probability model they
define. Given the calibration, the evidence estimate is unbiased, on every design. If, in addition, the rows are Gaussian, the signals fixed, every true predictor visible, and no
null predictor correlated with a true one beyond order $n^{-1/2}$, even given $15$ others, the printed estimates are
consistent: the errors of the evidence estimate (relative and log error) and of the inclusion probabilities, multiplied
by the indicator that the method uses the sampler and its calibration succeeds, are of order
$n^{-1/2}e^{C_\rho\sqrt{\log n}}=n^{-1/2+o(1)}$; the probability that this event occurs and the median model is wrong tends
to zero, and so does the probability that the run completes and the reported model is wrong. Nothing is claimed about how
likely the event is, or conditionally on it.}

\why{The numbers a user reads come with guarantees: the evidence estimate is right on average whatever the design, and
on designs without leakage, on the event that the calibration succeeds, its relative error tends to zero, the
inclusion probabilities tend to $0$ or $1$ and the median model equals the true model with probability tending to one,
and a completed run reports the true model with probability tending to one (no claim is made about how likely that
event is).}

\subsection{Unbiased evidence, on every design}\label{sec:unbiased}

\begin{proposition}[Each island's evidence estimate is unbiased]\label{prop:unbiased}
Under Convention~\as{C3}, given the calibration, each island's evidence estimate $\Zhat_r$ is unbiased for $Z$, for every
design and data set.
\end{proposition}

\asmlist{Convention \as{C3} only; no condition on the design (Appendix~\ref{app:conventions}).}

\status{\textbf{Proved and independently checked}. Registry: Appendix~\ref{app:status}, row 3.5.}

This is the classical unbiasedness of sequential Monte Carlo estimates of normalising constants \citep{DelMoral2004},
established for the method as coded, calibration included: the calibration is frozen before the final run draws its
random numbers. By linearity the equal-weight pooled estimate is then unbiased too: on $\Acal$,
$\E[\Zhat\mid\text{calibration}]=Z$.

\status{Pooled form: \textbf{proved and independently checked}; the frozen-run identities behind both forms are
proved and independently checked at the code of record. Registry: Appendix~\ref{app:status}, row 3.5.}

\subsection{Consistent estimates, on designs without leakage}\label{sec:outputs}

\paragraph{Scope.} The next theorem needs \as{NL} below: no null predictor leaks signal. Of our designs only the
independent ones (S1 and IND) satisfy it. On our block and AR designs Result~3 and Proposition~\ref{prop:unbiased} hold,
and no theorem yet covers the consistency of the printed inclusion probabilities, median model and evidence estimate; on the factor designs
only Results~1--2 and Proposition~\ref{prop:unbiased} hold (\SA{3.4}).

\begin{assumption}{NL}{No leakage, as seen by the start law}
Every null predictor $l\notin\Sz$ has a deterministic \emph{neighbourhood} $\Hl_l\supseteq\{l\}\cup\Sz$ with
$\abs{\Hl_l}\le e^{\kappa_H\sqrt{\log n}}$ for a fixed $\kappa_H$, such that, for every $U\not\ni l$ with $\abs{U}\le15$ and
$U_0:=U\cap\Hl_l$: (i) \emph{no leak:} $\abs{\Cov(W_l,W_j\mid W_{U_0})}\le C_Ln^{-1/2}$ for every $j\in\Sz\setminus U_0$,
with a fixed $C_L$; and (ii) \emph{locality:} $\abs{\Cov(W_l,W_v\mid W_{U_0})}\le r^\circ_n$ and
$\abs{\Cov(W_j,W_v\mid W_{U_0})}\le r^\circ_n$ for every $v\in U\setminus\Hl_l$ and $j\in\Sz\setminus U_0$, where
$r^\circ_n(n\log n)^{1/4}\to0$. \hfill [\SA{3.4}]
\end{assumption}

The theorem concerns runs that use the sampler; enumeration is not covered. Every completed run lies in the event
$\Acal$ of Section~\ref{sec:method}. On $\Acal$, $\Zhat$, $\PIPhat$ and
$\MPMhat$ denote the pooled values the four islands would produce if the run did not stop; on completed runs they
are the printed values. Let $\rho_n:=n^{-1/2}e^{C_\rho\sqrt{\log n}}$, with $C_\rho$ depending only on the fixed quantities of the
assumptions.

\begin{theorem}[The printed estimates are consistent]\label{thm:outputs}
Assume \as{GR}, \as{DIM} with $p>\sz$, \as{FSt}, \as{EFS}, \as{V}, \as{NL}, and Conventions~\as{C1} and~\as{C3}. Then,
under the joint law of the design, the noise and the method's random numbers,
\begin{gather}
  \ind{\Acal}\Big(\big|\Zhat/Z-1\big|+\big|\log\Zhat-\log Z\big|+\max_{j\in\set{p}}\big|\PIPhat_j-\ind{j\in\Sz}\big|\Big)
    =\OP(\rho_n),\label{eq:outputs-rate}\\
  \P\big(\Acal\cap\{\MPMhat\neq\Sz\}\big)\to0,\qquad \P\big(\text{the run completes and }\Shat\neq\Sz\big)\to0,
    \label{eq:outputs-sel}
\end{gather}
and~\eqref{eq:outputs-rate} also holds for the printed values multiplied by the indicator of a completed run.
\end{theorem}

\asmlist{(1)~\as{GR}, \as{FSt}, \as{EFS}, \as{V} as in Result~3; (2)~\as{DIM} [\SA{3.1}]; (3)~\as{NL} no leakage
[\SA{3.4}]; (4)~Conventions \as{C1}, \as{C3} (Appendix~\ref{app:conventions}).}

\status{\textbf{Proved and checked, conditional on} \samplerconds. Registry: Appendix~\ref{app:status}, row 3.6.}

\begin{remark}[What is not claimed]\label{rem:outputs}
(i) No lower bound on $\P(\Acal)$ is asserted: the statements are about intersections with $\Acal$ or with completion, not
conditional on them.
(ii) The rate $\rho_n$ is the start law's rate; it is slower than the posterior rate $\rn$, by the factor
$e^{C_\rho\sqrt{\log n}}/\sqrt{\log n}$.
(iii) Pooling cannot make things worse: the pooled inclusion probabilities are never farther from the exact ones, in the
largest coordinate, than those of the worst island (proved and checked, conditional on the same two sampler results as
Theorem~\ref{thm:outputs}; registry row 4.1).
\end{remark}


\bibliography{../shared/refs}

\appendix
% =====================================================================================================
\section{Supplement: further results}\label{sec:supp}

Each result below comes with one line on why it matters. Statements are as in the first version of this section; the
status notes give their standing.

\subsection{The posterior}\label{sec:supp-post}

\why{Without the beta-min condition, the posterior still rejects over-fitted models and still keeps every true predictor
that is strong enough; specificity, however, can fail.}

\begin{remark}[Without beta-min]\label{rem:no-betamin}
Under the other assumptions of Theorem~\ref{thm:post-fixed}, the posterior mass of models that strictly contain $\Sz$ is
$\OP(\rn)$, and every true predictor with $\theta_j^{*2}\ge c\,\sigs^2\log n/(n\varphi_0^2)$, $c>c_0(\sz,A)$, has
$1-\PIP_j=\OP(\rn)$. Specificity can fail without beta-min: a numerical counterexample, reported with the proof, shows this.
\end{remark}

\status{Part of the posterior concentration result, with the status of the in-probability form of
Theorem~\ref{thm:post-fixed}.}

\why{It replaces the eigenvalue floor away from the true model, and the beta-min condition, by a floor on the true model
alone and a margin on the unexplained signal; it keeps the vanishing floor \as{NC}.}

\begin{localassumption}[Floor on the true model and an identifiability margin]\label{asm:margin}
\begin{asmitems}
  \item $\lmin(\Rhat_{\Sz})\ge\varphi_0^2$ for a fixed $\varphi_0\in(0,1]$.
  \item Every model $S\in\Om$ that misses a true predictor and has at most $j_0:=\lceil(2A+1)\sz\rceil+1$ predictors
    outside $\Sz$ leaves unexplained signal $\norm{(I-P_S)\mu}^2\ge c_{\mathrm{id}}\,\sigs^2\log n$, for a fixed
    $c_{\mathrm{id}}>c_1(\sz,A)$, an explicit constant (Appendix~\ref{app:constants}).
\end{asmitems}
\end{localassumption}

\begin{theorem}[Posterior concentration under an identifiability margin]\label{thm:post-margin}
Let the design be fixed and assume \as{GL}, \as{DIM}, \as{SP}, \as{BC}, \as{NC} and Assumption~\ref{asm:margin}. Then
the conclusions~\eqref{eq:post-op}--\eqref{eq:post-exp} of Theorem~\ref{thm:post-fixed} hold, the expectation bound with
constants depending only on $(\sz,A,\sigs^2,\bar\theta,\varphi_0,g_0,c_{\mathrm{id}})$.
\end{theorem}

\status{\textbf{Proved and checked, conditional on the same two supporting results} (the lower-bound proposition and the
size-tail lemma, not yet independently checked) as the expectation bound of Theorem~\ref{thm:post-fixed}.
\bridgepending{} Registry: Appendix~\ref{app:status}, row 1.1.}

Under \as{EF} and \as{BM}(a), Assumption~\ref{asm:margin} holds with $c_{\mathrm{id}}=C_\beta^2\varphi_0^2(1-1/n)$. So
Theorem~\ref{thm:post-margin} contains Theorem~\ref{thm:post-fixed} when $C_\beta^2\varphi_0^2>c_1$; when
$c_0<C_\beta^2\varphi_0^2\le c_1$, Theorem~\ref{thm:post-fixed} applies and Theorem~\ref{thm:post-margin} is not known to
apply: the two coexist, and neither contains the other.

\why{It says what the posterior selects when some true predictors are missing from the data: the predictors used by the
best linear predictor built from the available ones.}

In this result the method is run on the response $\yr$ and a set $\mathcal C$ of $p$ candidate predictors, with rows
$x_{\mathcal C,i}$ collected in $x_{\mathcal C}$; $\Sigma$ and $R$ are the covariance and the correlation matrix of a
candidate row.

\begin{localassumption}[Candidate-only fitting]\label{asm:omitted}
\begin{asmitems}
  \item \emph{Model.} The candidate rows are i.i.d.\ $\Normal_p(m_{\mathcal C},\Sigma)$, and, given them, the responses are
    independent with $\yr_i\mid x_{\mathcal C}\sim\Normal(\alpha'+x_{\mathcal C,i}^\top\beta^{\mathcal C},\sigma^2_{\mathcal C})$,
    $\sigma^2_{\mathcal C}>0$. We write $\Sz^{\mathcal C}:=\supp\beta^{\mathcal C}$ and $s_{\mathcal C}:=\abs{\Sz^{\mathcal C}}$.
  \item \emph{Floor.} $\lmin(R_T)\ge c_{\mathcal C}$ for every $T\subseteq\mathcal C$ with
    $\abs{T}\le3(A+1)(s_{\mathcal C}+1)\log n$, for a fixed $c_{\mathcal C}>0$.
  \item \emph{Signal relative to the residual variance.} $\beta^{\mathcal C\top}\Sigma\beta^{\mathcal C}\le
    M_{\mathcal C}\sigma^2_{\mathcal C}$ for a fixed $M_{\mathcal C}$, and $n\min_{j\in\Sz^{\mathcal C}}\Sigma_{jj}(\beta^{\mathcal
    C}_j)^2/(\sigma^2_{\mathcal C}\log n)\to\infty$.
\end{asmitems}
\end{localassumption}

\begin{theorem}[Omitted variables]\label{thm:omitted}
Assume \as{DIM} (with $p$ the number of candidates), Assumption~\ref{asm:omitted}, and that $s_{\mathcal C}\ge1$ is fixed. No
condition is placed on how $m_{\mathcal C}$, $\Sigma$, $\beta^{\mathcal C}$ or $\sigma^2_{\mathcal C}$ vary with $n$. Then
\[
  \E\big[1-\pi(\Sz^{\mathcal C}\mid y)\big]=O\Big(\sqrt{\frac{\log n}{n}}\Big),\qquad
  \P\big(\HPM\neq\Sz^{\mathcal C}\big)=O\Big(\sqrt{\frac{\log n}{n}}\Big).
\]
\end{theorem}

\status{\textbf{Proved and checked, conditional on} the expectation bound of Theorem~\ref{thm:post-fixed} (hence on its two
supporting results) and on two lemmas on Gaussian designs that are verified and passed a council review. \bridgepending{}
Registry: Appendix~\ref{app:status}, row 1.4.}

For example, if candidates and omitted predictors are jointly Gaussian, $\beta^{\mathcal C}=\beta^*_{\mathcal C}+
\Sigma_{\mathcal C\mathcal C}^{-1}\Sigma_{\mathcal C\mathcal M}\beta^*_{\mathcal M}$: each candidate absorbs the part of the
omitted signal it can predict, and the rest is added to the noise variance $\sigma^2_{\mathcal C}$.

\subsection{The distribution}\label{sec:supp-dist}

\why{It shows what the collinearity condition of Result~2 buys: dropping it costs one logarithmic factor in the radius,
and the posterior also puts almost no mass on large models.}

\begin{proposition}[Level-$M$ form, and the case without the collinearity condition]\label{prop:dist}
\textup{(i)} There is a constant $C_{\mathrm s}$, depending only on $(s,A,L,c_\psi,\bar b)$, such that the following holds. The
constants of Theorem~\ref{thm:dist} serve every level $M$ with $M_0\le M\le\sqrt n/\log n$, and $M$ may depend on $n$:
whenever $n\ge n_0$ and Assumption~\ref{asm:dist} holds with the window $q_n$ of this $M$,~\eqref{eq:dist-i} holds, and also
\begin{equation}\label{eq:dist-ii}
  \P^*\Big[\,\Pi\big(\abs{S}>s+C_{\mathrm s}M\ \big|\ \Dn\big)>n^{-c_{\mathrm H} M}\Big]\ \le\ n^{-c_{\mathrm H} M}.
\end{equation}
The two exceptional sets can be taken to be the same, and the bounds are uniform over all laws of the data that satisfy
Assumption~\ref{asm:dist} with the given constants.
\textup{(ii)} If Assumption~\ref{asm:dist}\ref{asm:dist-nec} is dropped and Assumption~\ref{asm:dist}\ref{asm:dist-subg} is
required only for single predictors, then~\eqref{eq:dist-i} holds with radius $\sqrt{M\log n\cdot\log(\praw\vee n)/n}$,
and~\eqref{eq:dist-ii} with $s+C_{\mathrm s}M\log n$ in place of $s+C_{\mathrm s}M$.
\end{proposition}

\status{\textbf{Proved and independently checked} (part of the same result as Theorem~\ref{thm:dist}). Registry:
Appendix~\ref{app:status}, row 2.1.}

\why{With a fixed eigenvalue floor and every true signal above the detection scale, the distribution is learned at the
parametric rate.}

\begin{theorem}[Parametric rate when the true predictors are detectable]\label{thm:dist-param}
Assume Assumption~\ref{asm:dist} with the window $q_n$ of level $M=1$, with exactly $s\ge1$ nonzero coordinates of $\bstar$;
that $\lmin(\Corr(x_T))\ge\psi$ for a fixed $\psi>0$ and every set $T$ of at most $(2A+3)s+2$ predictors; and that
$\min\{\abs{\beta^*_j}\sd(x_j)/\sigs:\beta^*_j\neq0\}\ge b_n$ with $b_n^2n/\log n\to\infty$. Then
$\Pi(d_H(f,f^*)>M_n/\sqrt n\mid\Dn)\to0$ in $\P^*$-probability, for every sequence $M_n\to\infty$.
\end{theorem}

\status{\textbf{Proved and checked, conditional on} Theorem~\ref{thm:post-fixed} (in-probability form), which is proved
and independently verified. Registry: Appendix~\ref{app:status}, row 2.1.}

\why{The distribution guarantee carries over to the list of models that a completed run returns, whenever that list
holds a non-negligible share of the posterior.}

A completed run returns a list $\mathcal T$ of at most $40p+1$ distinct models with their log scores; let
$\pi(\mathcal T\mid\Dn)$ be the posterior probability of the list, and $\hat\Pi$ the posterior restricted to it (the models
of $\mathcal T$ weighted in proportion to their exponentiated log scores, each with its exact conditional posterior of
$(\beta,\sigma^2)$; $\hat\Pi$ is not computed by the method).

\begin{corollary}[The output of the sampler]\label{cor:dist-output}
Assume the conditions of Proposition~\ref{prop:dist}(i), and that the returned log scores are computed in exact
arithmetic. With $\P^*$-probability at least $1-n^{-c_{\mathrm H}M}$, every completed run, whatever its random numbers, with
$\pi(\mathcal T\mid\Dn)\ge n^{-c_{\mathrm H}M/2}$ satisfies $\hat\Pi(d_H(f,f^*)>\sqrt{M\log n/n})\le n^{-c_{\mathrm H}M/2}$. If a
completed run scored $\Sz$ and $\pi(\Sz\mid\Dn)>\tfrac12$, then the list contains $\Sz$, $\Shat=\Sz$, and $\pi(\mathcal
T\mid\Dn)>\tfrac12$.
\end{corollary}

\status{\textbf{Proved and checked, under exact arithmetic}; the description of the output was checked against the code.
Registry: Appendix~\ref{app:status}, row 2.1.}

Whether the list carries enough posterior probability is a property of the run; the corollary does not assert it.

\subsection{The method}\label{sec:supp-method}

\why{It explains why the sampler starts in the right place: without leakage, the start law already concentrates on the
true model.}

\begin{theorem}[The start law concentrates]\label{thm:start}
Assume \as{GR}, \as{DIM} with $p>\sz$, \as{FSt}, \as{V}, \as{NL}, and $\lmin(\Sigma_{n,S})\ge\varphi_{\mathrm b}^2$ for every
$S$ with $\abs{S}\le16+\sz$, for a fixed $\varphi_{\mathrm b}>0$. Then $\P(\pi_0(\Sz)\ge1-n^{-c})\to1$ for every
$c<\tfrac12$. If moreover $\Cov(W_l,W_j)=0$ for every null $l$ and true $j$, then $1-\pi_0(\Sz)=\OP(\rn)$. The exponent
$\tfrac12$ cannot be improved: $1-\pi_0(\Sz)\ge\tfrac13(n+1)^{-1/2}$ for every valid data set with $n\ge4$ and $\sz<K$.
\end{theorem}

\status{\textbf{Proved and independently checked}. Registry: Appendix~\ref{app:status}, row 3.3.}

With leakage, the start law concentrates instead on the models between the true model and the leakage set: if \as{NL}
holds with $\Sz$ replaced by a deterministic $Q\supseteq\Sz$ with $\abs{Q}\le K$, then
$\P\big(\pi_0(\Lat(Q))\ge1-(\abs{Q}+1)n^{-1/2}e^{C\sqrt{\log n}}\big)\to1$.

\why{It is the only reported-model statement that does not assume Gaussian rows: its conditions are on the sample
$(X,\tstar)$.}

For a set $U$ and a predictor $l\notin U$, the \emph{signal gain} $\delta^2_{U,l}:=\norm{(P_{U\cup\{l\}}-P_U)\mu}^2$ is the
increase in explained mean signal when $x_l$ joins $X_U$; the \emph{residual scale} is $\tau^2_U:=\sigs^2+\norm{(I-P_U)\mu}^2/(n-1)$;
and $\hat\beta_{U,l}:=(X_U^\top X_U)^{-1}X_U^\top x_l$, padded with zeros.

\begin{localassumption}[Fixed-design visibility and leakage]\label{asm:reported-fixed}
\begin{asmitems}
  \item\label{asm:rf-floor} \emph{Floors.} $\lmin(\Rhat_T)\ge\varphi_0^2$ for every $T$ with $\abs{T}\le16$ and for every
    $T=Q\cup\{l\}$, $l\in\set{p}$, with $Q$ as in~\ref{asm:rf-leak}.
  \item\label{asm:rf-vis} \emph{Visibility.} There is a deterministic $\omega_n\to\infty$ such that
    $\delta^2_{U,j}\ge\omega_n\log n$ for every $j\in\Sz$ and every $U\not\ni j$ with $\abs{U}\le15$.
  \item\label{asm:rf-leak} \emph{Bounded leakage.} There are a deterministic $Q\supseteq\Sz$ (with $\abs{Q}\le K$ and
    ($Q=\set{p}$ or $K\ge2\abs{Q}+32\sz+1$) for large $n$) and a fixed $\eta\in[0,\eta_A)$,
    $\eta_A:=(\sqrt{2A+1}-\sqrt{2A})^2$, such that $\delta^2_{U,l}\le\eta\,\tau^2_U\log n$ for every $l\notin Q$ and every
    $U\not\ni l$ with $\abs{U}\le15$.
  \item\label{asm:rf-size} \emph{Size of the leakage set.} One of: (d1) $\abs{Q\setminus\Sz}\le(1-\delta_Q)\log n$ for a
    fixed $\delta_Q\in(0,1)$; or (d2) $\abs{Q}\le n^b$ with $b<\tfrac12-\tfrac\eta2-\sqrt{2A\eta}$,
    $\theta^{*2}_{\min}n/(\abs{Q}+\log n)\to\infty$, and local regressions: for every $l$ a deterministic set
    $N_l\subseteq Q\setminus\{l\}$ with $\abs{N_l}\le h\log_2n$, for a fixed $h$ with $h+b<\tfrac12$, and a deterministic $L_n\to0$ with
    $L_n^2\log(\abs{Q}+1)=o(\log n)$, such that $\norm{\hat\beta_{T,l}-\hat\beta_{T\cap N_l,l}}_1\le L_n$ for every
    $T\in\Lat(Q)$ with $\abs{T}<K$ and $l\notin T$.
\end{asmitems}
\end{localassumption}

\begin{theorem}[The reported model equals the true model: fixed design]\label{thm:reported-fixed}
Let the design be fixed, with no raw predictor constant in the sample (an input condition of the delivery result; then
$p=\praw$). Assume \as{GL}, \as{DIM}, \as{SP}, \as{BC},
\as{EF}, \as{NC}, \as{BM}(b), Assumption~\ref{asm:reported-fixed} and Convention~\as{C1}. Then the conclusions of
Theorem~\ref{thm:reported} hold: there are data events $\mathcal G_n$ with $\P(\mathcal G_n)\to1$ on which every
admissible run that produces output reports $\Shat=\Sz$, for every seed and calibration outcome; and the delivery
clauses hold under \as{C2} on $\Dom$ and under \as{C2p} for every admissible $n\ge3$.
\end{theorem}

\status{\textbf{Proved and independently checked}. \bridgepending{} A revision of this statement, from a strengthening
that has passed its review, is pending. Registry: Appendix~\ref{app:status}, row 3.2.}

\why{A certified small model rarely misses a true predictor, under much weaker conditions than Result~3; it is the
guarantee about the coded method that remains on pervasive-factor designs, which Result~3 excludes.}

\begin{theorem}[Certified screening]\label{thm:screen}
Let the design be fixed and assume \as{GL} and Convention~\as{C1}; \as{SP} is not assumed. Let $k_n\in\{0,\dots,K-1\}$ be
a model-size budget, and assume $\lmin(\Rhat_S)\ge\kappa$ for every $S$ with $\abs{S}\le k_n+\sz$, $\tmin\ge\tlow>0$ and
$\norm{\tstar}_1\le M_1$, with $\kappa$, $\tlow$, $M_1$, $\sigs^2$ fixed, and $(k_n+1)\log p+\log n\le c_{\mathrm{SB}}\,n$, where
$c_{\mathrm{SB}}:=\kappa^2\tlow^2/\{200(M_1^2+\sigs^2)\}$. Then
\[
  \P\big(\text{the run produces output},\ \Shat\text{ is certified},\ \abs{\Shat}\le k_n,\ \Sz\not\subseteq\Shat\big)\to0 .
\]
\end{theorem}

\status{\textbf{Proved and independently checked}. Registry: Appendix~\ref{app:status}, row 3.7.}

No visibility, leakage, start or block condition is used. A second form replaces the signal bounds by a margin on the
noiseless gain of a missing true predictor, with no condition on $p$, $\sz$ or the null predictors (also proved and
independently checked). The bound is on an intersection, not conditional on certification, and a run is not claimed to certify a small model
with probability tending to one. In existing numerics (not a theorem) every completed run certified its reported model:
2,394 of 2,394 in the 2,400-fit study, 376 of 376 in the strong-factor study, 71 of 71 in a screening study. On the
strong-factor designs the certified model contained every true predictor in 61 of 122 fits at $n=50$, 18 of 20 at $n=75$
and 30 of 30 at $n=100$; at the constants of these designs the size budget of Theorem~\ref{thm:screen} requires
$n\gtrsim1{,}450$.

\why{It is the only recorded route to the printed inclusion probabilities and median model that does not need \as{NL};
its global-overlap premise is not yet derived from conditions on the design.}

\begin{remark}[An alternative route for the estimates]\label{rem:alt-route}
An argument that does not go through the evidence estimate gives the conclusions of Theorem~\ref{thm:outputs} for the
inclusion probabilities and the median and highest posterior models under a global-overlap condition (stated in full with
the proof).
\end{remark}

\status{\textbf{Proved and independently checked}; its premise is open. Registry: Appendix~\ref{app:status}, row 3.6.}

\subsection{Open questions}\label{sec:open}

None of these is needed for any statement above.
\begin{enumerate}[leftmargin=1.5em,itemsep=1pt]
  \item Remove the vanishing floor \as{NC} from Theorem~\ref{thm:post-margin}, and weaken its floor on the true model.
  \item Prove $\E[1-\pis{\Sz}]\gtrsim\rn$ on every design with a floor on the sets $\Sz\cup\{j\}$.
  \item Whether Assumption~\ref{asm:dist}\ref{asm:dist-nec} is needed for the rate $\rn$ of Result~2.
  \item Leakage sets of size $n^b$ with $b\ge\tfrac12$, and raw-scale signals in Results~3 and~4.
  \item Consistency of the printed estimates on designs with leakage (AR and block designs).
  \item The parallel method: an inheritance statement for an output-identical engine, and results for many islands and
    island weighting at very large $p$.
\end{enumerate}

% =====================================================================================================
\section{Notation}\label{app:notation}

\begin{longtable}{@{}p{0.28\linewidth}p{0.5\linewidth}l@{}}
\toprule
Symbol & Meaning & Where \\ \midrule \endhead
$n$; $p$; $\praw$ & sample size; retained (non-constant) predictors; raw predictors & \S\ref{sec:model} \\
$\alpha^*,\bstar,\sigs^2$; $\Sz$, $\sz$ & true intercept, coefficients, noise variance; true model and its size & \S\ref{sec:model} \\
$\Sig$, $W$ & population correlation matrix of a row; a centred population row & \S\ref{sec:model} \\
$X$, $y$, $\tstar$, $\mu$, $\tmin$ & standardised design and response; standardised coefficients; mean signal; smallest true $\abs{\theta^*_j}$ & \S\ref{sec:model} \\
$X_S$, $G_S$, $\Rhat_S$, $P_S$ & columns in $S$; Gram matrix; sample correlation matrix; projection & \S\ref{sec:model} \\
$\Om$, $K$ & model space; size cap $\min(p,\lfloor n/2\rfloor)$ & \S\ref{sec:model} \\
$\pis{S}$, $F(S)$, $Z$ & posterior over models; log score; evidence & \S\ref{sec:model} \\
$\PIP_j$, $\HPM$, $\MPM$ & inclusion probability; highest posterior model; median probability model & \S\ref{sec:model} \\
$\pi_0$, $\tSz$, $\Shat$ & start law; start mode; reported model & \S\ref{sec:method} \\
$\PIPhat_j$, $\Zhat$, $\MPMhat$ & pooled estimates and estimated median probability model & \S\ref{sec:method} \\
$\varphi_0,g_0,\bar\theta,C_\beta$; $\varphi$; $b_{\mathrm{lo}},b_{\mathrm{hi}}$ & constants of the fixed-design assumptions; population floor; bounds on raw coefficients & \S\ref{sec:standing}, \S\ref{sec:model-result} \\
$c_0(\sz,A)$, $c_1(\sz,A)$, $C_{\mathrm E}$ & beta-min and margin constants; rate constant & App.~\ref{app:constants} \\
$\Dn$, $\Pi(\cdot\mid\Dn)$, $\P^*$, $\zeta$, $d_H$ & data set; posterior of $(S,\beta,\sigma^2)$; law of the data; standardised predictors; Hellinger distance & \S\ref{sec:distribution} \\
$M$, $q_n$, $c_{\mathrm H}$, $M_0$ & level, window and constants of the distribution theorem & \S\ref{sec:distribution} \\
$\Cov(\cdot,\cdot\mid W_U)$, $\bar\beta_{U,l}$ & partial covariance; population regression coefficient & \S\ref{sec:reported} \\
$Q$, $\Lat(Q)$, $\Hl_l$; $\theta_{\mathrm{lo}},\theta_{\mathrm{hi}}$ & leakage set; its lattice; locality sets; bounds on standardised coefficients & \S\ref{sec:reported} \\
$\mathcal G_n$; $\Dom$ & data event of Result~3; delivery domain & \S\ref{sec:reported}; App.~\ref{app:conventions} \\
$\Acal$, $\rho_n$ & sampler used and calibration succeeds; rate $n^{-1/2}e^{C_\rho\sqrt{\log n}}$ & \S\ref{sec:outputs} \\
\bottomrule
\end{longtable}

\section{Conventions in full}\label{app:conventions}

The code has two fixed integer limits on the local-search move caps; \as{C2p} lifts them. They are not reached on the
domain
\[
  \Dom:=\big\{(n,\praw):\lceil\log_2(n+\praw+1)\rceil\le10^6,\ \ \lceil\log_2(n+\praw+1)\rceil+\lfloor n/2\rfloor\le2^{31}-1\big\}.
\]
The limits are a ceiling of $10^6$ on the base move cap $\ell_n:=\lceil\log_2(n+p+1)\rceil$, and the 32-bit integer bound
$2^{31}-1$ on the move cap $\ell_n+\abs{\tSz}$ of the search from the start mode. The first can bind only when $n+p+1$
exceeds $2^{10^6}$; the second only when the start mode has more than about $2.1\times10^9$ predictors, which needs
$n\gtrsim4.3\times10^9$, since $\abs{\tSz}\le\lfloor n/2\rfloor$. In the default configuration neither can be reached on
any input R can represent ($n,\praw\le2^{31}-1$), so $\Dom$ contains every such input; under \as{C2} without the lift,
the delivery clauses hold on $\Dom$ (Theorem~\ref{thm:reported-fixed} states this form).

\paragraph{Convention \as{C2}: delivery.} Used only by the clauses ``the run always produces output'' of
Theorems~\ref{thm:reported} and~\ref{thm:reported-fixed}. They are operational assumptions, not consequences of any
statistical condition.
\begin{enumerate}[label=(\roman*),leftmargin=2em]
  \item \emph{Idealised machine:} exact arithmetic at unit cost (including $\exp$ and $\log$); no integer, index,
    random-stream identifier or allocation overflow; wall time at most $c_{\mathrm t}\cdot(\text{operation count})+c_{\mathrm t0}$ and peak memory at
    most $c_{\mathrm m}\cdot(\text{words held})+c_{\mathrm m0}$ for fixed constants $(c_{\mathrm t},c_{\mathrm t0},c_{\mathrm m},c_{\mathrm m0})$; every allocation within
    the declared memory succeeds; the build and its dependencies work; no external interruption or worker crash.
  \item \emph{Covered configuration:} the default core of the software, with the local search from the start mode (cap
    $\ell_n+\abs{\tSz}$) and the rule that resolves ties in the start mode, implemented with
    $O\big(p\sum_b(2^{\abs{C_b}}+(K+1)^2)\big)$ operations and $O\big(\sum_b2^{\abs{C_b}}+J(K+1)\big)$ storage; the local
    search from the start mode starting at the smallest, lexicographically first maximiser of the start law; the default
    list of summaries only; certificate computations with declared polynomial operation and memory budgets, included in
    the allocation; serial execution or a fixed number of concurrent workers (memory envelope multiplied by that number);
    a fixed number of islands; the fallback temperature sequence switched on (the default).
  \item \emph{Random numbers and stops:} ideal independent uniform random numbers; the method's stop-and-return rule and
    stop-only resource guard as coded (a stopped run returns no estimates).
  \item \emph{Inputs and allocation:} a valid seed (an integer in $[1,2^{31}-1]$) and a declared memory and time allocation
    at least the method's explicit polynomial work and memory envelopes in $(n,\praw)$, including the memory checks made before the run and
    the bounds on the records and time of the local searches.
\end{enumerate}
The code has two fixed integer limits: a run is rejected when the base cap $\ell_n$ exceeds a fixed ceiling, or when the
move cap $\ell_n+\abs{\tSz}$ of the search from the start mode exceeds the integer range. Theorem~\ref{thm:reported-fixed} states its delivery clause under (i)--(iv) on the domain $\Dom$ where neither limit
is reached; Theorem~\ref{thm:reported} states it under \as{C2p}. \emph{Convention \as{C2p}} is the idealised machine of (i)
with these two limits lifted, everything else (other validators, the cap on the enumeration used to find the start mode, and its fallback, every
method cap) kept as coded; under it the delivery clauses hold for every $n$. With the limits in place, (i)--(iv) give delivery only on $\Dom$, not for every
$n$.

\paragraph{Convention \as{C3}: sampler.} Used only by Theorem~\ref{thm:outputs} and Proposition~\ref{prop:unbiased}.
\begin{enumerate}[label=(\roman*),leftmargin=2em,itemsep=1pt]
  \item The computer does exact arithmetic.
  \item Within each random stream, the random numbers are fresh, independent, ideal uniforms; the final run's stream for
    each island is independent of the data and of the calibration's streams.
  \item The coded start sampler returns exact draws from $\pi_0$, and the reported start normaliser equals the exact
    normalising constant of the unnormalised weights of $\pi_0$.
  \item The coded Markov moves realise their stated kernels, applied in a uniformly random order that is independent of
    the state.
  \item The outcome of the calibration is a function of the data and of the calibration's random streams, and each island's
    record is a measurable function of that outcome and of the island's stream.
  \item Statements about estimates refer to the potential (uninterrupted) run, which is identified with the reported
    values on completed runs.
\end{enumerate}

\section{Explicit constants}\label{app:constants}

The constants are copied from the statements of the underlying results; none is optimised. With
$j_0:=\lceil(2A+1)\sz\rceil+1$ and $C_x:=Aj_0+\sz+\log(j_0+1)+2$:
\begin{itemize}[leftmargin=1.5em]
  \item \emph{Beta-min constant:} $c_0(\sz,A)=\max\{4C_*,\,32[\sz(A+1)+2]\}$, where
    $C_*:=2\sz+3[Aj_0+\sz+\log(j_0+1)+1]=3C_x+2\sz-3$ is the noise constant of the proof: with probability at least
    $1-1/n$, every model that misses a true predictor and has at most $j_0$ extra predictors has noise term at most
    $C_*\sigs^2\log n$.
  \item \emph{Margin constant:} with $C_A:=8+3A$,
    \[
      c_1(\sz,A)=\Big(\sqrt{2C_x}+\sqrt{2C_x+C_Aj_0+(2/\log2)\{\sz(A+1)+2\}}\Big)^2 .
    \]
\end{itemize}

\ifshowstatus
\section{Status ledger (internal; removed for submission)}\label{app:status}

Registry: \texttt{claims.tsv} at tag \texttt{theory-v2026-11-freeze} (d40c9780; claims card v1.13). Card rows use the
reader numbering of the registry (Theorem~1 = Theorem~\ref{thm:post-fixed}, Theorem~2 = Theorem~\ref{thm:post-random},
Theorem~3 = Remark~\ref{rem:sharp}, Theorem~9 = Theorem~\ref{thm:dist}, Theorem~6 = Theorem~\ref{thm:reported},
Theorem~10 = Theorem~\ref{thm:outputs}). Long registry strings are cut with an ellipsis.

\IfFileExists{status_ledger_freeze.tex}{\begingroup\footnotesize\sloppy
\begin{longtable}{@{}p{0.15\linewidth}p{0.05\linewidth}p{0.52\linewidth}p{0.2\linewidth}@{}}
\toprule Result & Row & Registry status, verbatim (\texttt{claims.tsv}, tag \texttt{theory-v2026-11-freeze}) & Label \\ \midrule \endhead
Theorem 3.1 (random design) & 1.2 & \texttt{W3-POST:Thm-1@1.0.1}: "PROVED-CHECKED; council C3, MANIFEST 420c1cb0\ldots{}, astra=8.8 PASS, opus=8.7 PASS"; \texttt{W2-APP:POST@1.0.7}: "PROVED-CHECKED"; \texttt{W3-READER:R1@1.0.2}: "PROVED-CHECKED; council R3, MANIFEST 7aebdf2e\ldots{}, sol=9.5 PASS, opus=9.4 PASS" & proved and independently checked \\
Theorem 3.2, in-probability form & 1.1 & \texttt{SOP:TH09@1.0.2}: "PROVED-CHECKED (T03$\sharp$ only: the conjunction of W3-T03ACC:T03S.1--T03S.4, clause map CHECKPOINT\_4.md §2; council R1, MANIFEST 9714aea4\ldots{}, astra=9.5 PASS, opus=9.4 PASS; LC-A, LC-B\ldots{}"; \texttt{W3-POST:Thm-1@1.0.1}: "PROVED-CHECKED; council C3, MANIFEST 420c1cb0\ldots{}, astra=8.8 PASS, opus=8.7 PASS"; \texttt{W3-READER:R1@1.0.2}: "PROVED-CHECKED; council R3, MANIFEST 7aebdf2e\ldots{}, sol=9.5 PASS, opus=9.4 PASS" & proved and independently checked \\
Theorem 3.2, expectation bound & 1.1 & \texttt{W3-POST2:C@1.0.6}: "PROVED-CHECKED, conditional on W3-POST:Prop-LB, W2-APP:POST"; \texttt{W3-POST:Prop-LB@1.0.0}: "PROVED (by this seat, unchecked)"; \texttt{W3-POST:Lemma-1a'@1.0.1}: "PROVED-CHECKED; council C3, MANIFEST 420c1cb0\ldots{}, astra=8.8 PASS, opus=8.7 PASS"; \texttt{W2-APP:POST@1.0.7}: "PROVED-CHECKED" & proved and checked, conditional on one supporting result (Prop-LB) not yet independently checked \\
Remark 3.3 (Result 5) & 1.3 & \texttt{W3-PSHARP:A@1.0.0}: "PROVED-CHECKED; council R2, MANIFEST 4513e05a\ldots{}, sol=9.2 PASS, opus=9.3 PASS; correct:PASS(reviews/R2\_council\_sol.md, reviews/R2\_council\_opus.md, per-claim addenda: ACCEPT)"; \texttt{W3-PSHARP:B@1.0.0}: "PROVED-CHECKED; council R2, MANIFEST 4513e05a\ldots{}, sol=9.2 PASS, opus=9.3 PASS; correct:PASS(reviews/R2\_council\_sol.md, reviews/R2\_council\_opus.md, per-claim addenda: ACCEPT)"; \texttt{W3-PSHARP-R:AE@1.0.0}: "PROVED-CHECKED; council R2, MANIFEST fdfd4562\ldots{}, sol=9.5 PASS, opus=9.2 PASS"; \texttt{W3-PSHARP-R:A2C@1.0.0}: "PROVED-CHECKED; council R2, MANIFEST fdfd4562\ldots{}, sol=9.5 PASS, opus=9.2 PASS" & proved and independently checked (exact order: status of the expectation bound) \\
Theorem 4.2 & 2.1 & \texttt{W3-DIST:Thm-A@1.0.1}: "PROVED-CHECKED; council R2, MANIFEST 9719c817\ldots{} (= R2 c5cd9cdd re-frozen CITATION-ONLY as R2r), astra=9.3 PASS, opus=9.2 PASS" & proved and independently checked \\
Theorem 5.3 & 3.1 & \texttt{W3-RM2-R:RM2-R-plus@1.0.1}: "PROVED-CHECKED; council R3, MANIFEST 5d6d3c90\ldots{} (= R3 b06a8588 re-frozen as R3r), sol=9.5 PASS, opus=9.4 PASS"; \texttt{W3-RM2-R:RM0-plus@1.0.1}: "PROVED-CHECKED; council R3, MANIFEST 5d6d3c90\ldots{} (= R3 b06a8588 re-frozen as R3r), sol=9.5 PASS, opus=9.4 PASS"; \texttt{W3-RM2-R:DLV-0@1.0.1}: "PROVED-CHECKED; council R3, MANIFEST 5d6d3c90\ldots{} (= R3 b06a8588 re-frozen as R3r), sol=9.5 PASS, opus=9.4 PASS"; \texttt{C-DLV:L2-S3@1.0.1}: "CODE-CHECKED (second seat, code reading; holds at pin 4676af6 after amendments A-1--A-7); council R2, MANIFEST 12802ac7\ldots{}, sol=9.3 PASS, opus=9.4 PASS (= track A\_C C1, MANIFEST\_\ldots{}"; \texttt{C-DLV:P1-Inventory@1.1.1}: "CODE-CHECKED (second seat, code reading; holds at pin 4676af6 after amendments A-1--A-7); council R2, MANIFEST 12802ac7\ldots{}, sol=9.3 PASS, opus=9.4 PASS (= track A\_C C1, MANIFEST\_\ldots{}"; \texttt{W3-RM2:RM2-R@1.0.3}: "PROVED-CHECKED, conditional on W2-APP:POST; input versions W2-APP:POST@1.0.1; council R1, MANIFEST 00bbeab8\ldots{}, astra=9.5 PASS, opus=8.8 PASS; correct:PASS(reviews/C\_R1\_council\_\ldots{}"; \texttt{W2-APP:POST@1.0.7}: "PROVED-CHECKED"; \texttt{W3-READER:R2@1.0.1}: "PROVED-CHECKED; council R3, MANIFEST 7aebdf2e\ldots{}, sol=9.5 PASS, opus=9.4 PASS"; \texttt{W3-AUDIT-AX:REP@1.0.0}: "PROVED-CHECKED; council R2, MANIFEST 5db04d3a\ldots{}, sol=9.4 PASS, opus=9.5 PASS" & proved and independently checked; no open condition \\
Proposition 6.1 and the pooled form & 3.5 & \texttt{W1-N1:Thm-N1@1.0.0}: "PROVED-CHECKED"; \texttt{W3-READER:PoolZ@1.0.1}: "PROVED-CHECKED; council R3, MANIFEST 7aebdf2e\ldots{}, sol=9.5 PASS, opus=9.4 PASS"; \texttt{W3-AUDIT-AX:FR@1.0.0}: "PROVED-CHECKED; council R2, MANIFEST 5db04d3a\ldots{}, sol=9.4 PASS, opus=9.5 PASS" & proved and independently checked \\
Theorem 6.2 & 3.6 & \texttt{W3-SAMP:C@1.1.5}: "PROVED-CHECKED, conditional on W1-N1:B, W1-N1:F; council R1, MANIFEST 68480a72\ldots{}, astra=9.1 PASS, opus=8.8 PASS"; \texttt{W3-SAMP:TR@2.1.5}: "PROVED-CHECKED; council R1, MANIFEST 68480a72\ldots{}, astra=9.1 PASS, opus=8.8 PASS"; \texttt{W3-SAMP:LT\#@2.0.5}: "PROVED-CHECKED, conditional on W1-N1:B, W1-N1:F; council R1, MANIFEST 68480a72\ldots{}, astra=9.1 PASS, opus=8.8 PASS"; \texttt{W3-SAMP:D@2.1.5}: "PROVED-CHECKED, conditional on W1-N1:B; council R1, MANIFEST 68480a72\ldots{}, astra=9.1 PASS, opus=8.8 PASS"; \texttt{W3-START2-R:IF-START@1.0.0}: "PROVED-CHECKED; council R2, MANIFEST 12bb52cd\ldots{} (= R2 90886d06 re-frozen CITATION-ONLY as R2r), sol=9.0 PASS, opus=9.3 PASS"; \texttt{W2-APP:POST@1.0.7}: "PROVED-CHECKED"; \texttt{W3-READER:R4@1.0.1}: "PROVED-CHECKED, conditional on W1-N1:B, W1-N1:F; council R3, MANIFEST 7aebdf2e\ldots{}, sol=9.5 PASS, opus=9.4 PASS"; \texttt{W1-N1:B@1.0.0}: "PROVED under the named premises of STATEMENT §2 — cross-family checked (Codex gpt-6.1-sol: LC1, MP1 CONTINUE, SW1, council C1 9.3) and council-passed (Claude Opus 5.5: C1 8.9), round 1, 2026-11-08."; \texttt{W1-N1:F@1.0.0}: "PROVED under the named premises of STATEMENT §2 — cross-family checked (Codex gpt-6.1-sol: LC1, MP1 CONTINUE, SW1, council C1 9.3) and council-passed (Claude Opus 5.5: C1 8.9), round 1, 2026-11-08." & proved and checked, conditional on W1-N1:B, W1-N1:F (themselves proved under named premises and independently checked) \\
Remark 6.3(iii) & 4.1 & \texttt{W3-SAMP:LT\#@2.0.5}: "PROVED-CHECKED, conditional on W1-N1:B, W1-N1:F; council R1, MANIFEST 68480a72\ldots{}, astra=9.1 PASS, opus=8.8 PASS" & proved and checked, conditional on W1-N1:B, W1-N1:F \\
Theorem A.3 & 1.1 & \texttt{W3-POST2:E@1.0.6}: "PROVED-CHECKED, conditional on W2-APP:POST"; \texttt{W3-POST2:C@1.0.6}: "PROVED-CHECKED, conditional on W3-POST:Prop-LB, W2-APP:POST" & in probability: proved and independently checked; expectation bound: conditional on Prop-LB \\
Theorem A.5 & 1.4 & \texttt{W3-MISR-R:CorGb-plus@1.0.1}: "PROVED-CHECKED, conditional on W3-POST2:C, W3-POST2:E, W3-POST2:A1, W3-POST2:A0, C-MIS:GDb, C-MIS:Prop-P; council R3, MANIFEST 052ea6df\ldots{}, sol=10.0 PASS, opus=9.6 PASS"; \texttt{W3-MISR:MIS(i)@1.0.1}: "PROVED-CHECKED, conditional on W3-POST2:C; council R2, MANIFEST 217d6708\ldots{}, sol=9.2 PASS, opus=8.9 PASS; correct:PASS(reviews/R2\_council\_sol.md, reviews/R2\_council\_opus.md, per\ldots{}"; \texttt{C-MIS:GD(b)@1.0.0}: "R02 VERIFIED-WITH-AMENDMENTS (applied); council PASS"; \texttt{C-MIS:Prop-P@1.0.0}: "R02 VERIFIED; R04 wording fixes checked in R05; council PASS" & proved and checked, conditional on the expectation bound (hence Prop-LB) and two verified Gaussian-design lemmas \\
Proposition A.6 & 2.1 & \texttt{W3-DIST:Thm-A@1.0.1}: "PROVED-CHECKED; council R2, MANIFEST 9719c817\ldots{} (= R2 c5cd9cdd re-frozen CITATION-ONLY as R2r), astra=9.3 PASS, opus=9.2 PASS" & proved and independently checked \\
Theorem A.7 & 2.1 & \texttt{W3-DIST:Thm-B@1.0.1}: "PROVED-CHECKED; council R2, MANIFEST 9719c817\ldots{} (= R2 c5cd9cdd re-frozen CITATION-ONLY as R2r), astra=9.3 PASS, opus=9.2 PASS" & proved and independently checked \\
Corollary A.8 & 2.1 & \texttt{W3-DIST:Cor-C@1.0.1}: "PROVED-CHECKED, conditional on W1-N1:B, W1-N1:F, EX-A (exact arithmetic); council R2, MANIFEST 9719c817\ldots{} (= R2 c5cd9cdd re-frozen CITATION-ONLY as R2r), astra=9.3 PASS, opus=9\ldots{}" & proved and checked, conditional on W1-N1:B, W1-N1:F and exact arithmetic \\
Theorem A.9 & 3.3 & \texttt{W3-START2-R:START-R@1.0.0}: "PROVED-CHECKED; council R2, MANIFEST 12bb52cd\ldots{} (= R2 90886d06 re-frozen CITATION-ONLY as R2r), sol=9.0 PASS, opus=9.3 PASS"; \texttt{W3-START2-R:BR@1.0.0}: "PROVED-CHECKED; council R2, MANIFEST 12bb52cd\ldots{} (= R2 90886d06 re-frozen CITATION-ONLY as R2r), sol=9.0 PASS, opus=9.3 PASS"; \texttt{W3-START:FLOOR@1.0.0}: "PROVED-CHECKED" & proved and independently checked \\
Theorem A.11 & 3.2 & \texttt{W3-RM2-R:RM2-F-plus@1.0.1}: "PROVED-CHECKED; council R3, MANIFEST 5d6d3c90\ldots{} (= R3 b06a8588 re-frozen as R3r), sol=9.5 PASS, opus=9.4 PASS"; \texttt{W3-RM2:RM2-F@1.0.3}: "PROVED-CHECKED; council R1, MANIFEST 00bbeab8\ldots{}, astra=9.5 PASS, opus=8.8 PASS; correct:PASS(reviews/C\_R1\_council\_astra.md, reviews/C\_R1\_council\_opus.md, per-claim addenda)" & proved and independently checked; no open condition \\
Theorem A.12 & 3.7 & \texttt{W3-SCR:A1@1.0.0}: "PROVED-CHECKED"; \texttt{W3-SCR:A-inf@1.0.0}: "PROVED-CHECKED"; \texttt{W3-SCR:SCR-A@1.0.0}: "PROVED-CHECKED" & proved and independently checked \\
Remark A.13 & 3.6 & \texttt{W3-SAMP:TR@2.1.5}: "PROVED-CHECKED; council R1, MANIFEST 68480a72\ldots{}, astra=9.1 PASS, opus=8.8 PASS" & proved and independently checked; premise P-OVL$\infty$ open (track A T11) \\
\bottomrule
\end{longtable}
\endgroup
}{[Ledger table.]}

\paragraph{What changed at the freeze.} Since the first draft of this version (statuses of main a7fc5093): the
reader-form bridges are independently checked; the strengthened masters of Theorems~\ref{thm:reported}
and~\ref{thm:reported-fixed} are recorded, and the code-reading audits behind the delivery route are checked, so neither
theorem carries an open condition; the representability lemma is recorded, and the delivery clause of
Theorem~\ref{thm:reported} takes its one-sentence form; the size-tail lemma, the frozen-run identities and the pooled
evidence estimate are independently checked; Theorem~\ref{thm:dist-param} no longer carries a condition;
Theorem~\ref{thm:outputs}, Corollary~\ref{cor:dist-output} and the worst-island bound now carry two explicit conditions,
results on the coded sampler that are themselves proved under named premises and independently checked.

\paragraph{What remains after the freeze.}
\begin{enumerate}[leftmargin=1.5em,itemsep=1pt]
  \item The lower-bound proposition behind~\eqref{eq:post-exp} awaits its independent check.
  \item The superseded masters of Theorems~\ref{thm:reported} and~\ref{thm:reported-fixed} are retired in the registry
    (deferred past the freeze; no effect on the statements).
  \item The global-overlap premise of Remark~\ref{rem:alt-route} stays a named open premise.
  \item \emph{Planned, not decided:} if the store-restricted result (track D, \#5$'$) is proved, Theorem~\ref{thm:outputs}
    may become a remark, and consistency of the summaries computed from the returned list would be stated on every
    design of Result~3. Section~\ref{sec:estimates} is written so that this needs no other change.
\end{enumerate}
\fi


\end{document}
